% phil479-screen.tex

% driver file phil479-screen.tex to produce the Logic II textbook on
% with same type block as in printed version, but but with on-screen
% features (color, links, etc).

% We use the memoir class for maximal flexibility of layout, but any
% class will do

\documentclass[11pt,openany]{memoir}

\RequirePackage{amsthm}
\RequirePackage{xcolor}
\RequirePackage{mdframed}
\RequirePackage[full]{textcomp}

\RequirePackage[pdftex,breaklinks,bookmarks,bookmarksopen,bookmarksopenlevel=1,colorlinks,urlcolor=dark-gray,linkcolor=reflex-blue]{hyperref}

%\usepackage[nonumberlist,toc,style=index]{glossaries}
%\makeglossaries


% let's set the whole thing in Baskervald X, with Universalis ADF
% Standard for sans-serif, and spread the lines a bit to make the text
% more readable

\usepackage[sfdefault]{universalis}
\usepackage[osf]{Baskervaldx} % oldstyle figures
\usepackage[bigdelims,baskervaldx]{newtxmath} 
\usepackage[cal=boondoxo]{mathalfa}

% Make sure we have a copyright symbol

\def\copyright{\textcircled{C}}

\def\oljobname{phil479}

% next line needed for Adobe Acrobat to open the PDF.  may have
% somethign to do with transparency in the PNG graphics files?

\pdfminorversion=4

% load textpos for the cover

\usepackage[absolute,overlay]{textpos}

% set stock & paper size with narrow margins

\setstocksize{22cm}{17cm}

\settrimmedsize{\stockheight}{\stockwidth}{*}
\settrims{0pt}{0pt}

% let's calculate the line length for 65 characters in \normalfont

\setlxvchars

% set the size of the type block to golden ratio calculated width

\settypeblocksize{*}{1.05\lxvchars}{1.62}

% set spine and and edge margin

\setlrmargins{*}{*}{1}
\setulmargins{60pt}{*}{*}
\setheaderspaces{*}{*}{1}

\checkandfixthelayout

% Chapter style

\makeatletter

\definecolor{leadbeater}{RGB}{248,154,14}
\definecolor{dkleadbeater}{RGB}{104,65,6}
\definecolor{ltleadbeater}{RGB}{248,203,137}
\newlength{\barlength}
\makechapterstyle{leadbeater}{%
  \setlength{\afterchapskip}{40pt}
  \setlength{\beforechapskip}{50pt}
    \setlength{\midchapskip}{10pt}
  \renewcommand*{\afterchapternum}{\par\nobreak\vskip 0pt}
  \renewcommand*{\chapnamefont}{\fontsize{14pt}{0pt}\selectfont\bfseries\sffamily}
  \renewcommand*{\chapnumfont}{\fontsize{14pt}{0pt}\selectfont\bfseries\sffamily}
  \renewcommand*{\chaptitlefont}{\normalfont\fontsize{48pt}{48pt}\selectfont\bfseries\itshape\color{leadbeater}}
  \renewcommand*{\printchaptername}{%
    \chapnamefont\MakeUppercase{\@chapapp}}
  \renewcommand*{\printchaptertitle}[1]{%
    \chaptitlefont ##1\\[-\baselineskip]%
    \hspace*{-20pt}%
    \smash{\color{leadbeater}\rule{7pt}{300pt}}}
}

\renewcommand*{\partnamefont}{\fontsize{24pt}{0pt}\selectfont\bfseries\sffamily}
\renewcommand*{\partnumfont}{\fontsize{24pt}{0pt}\selectfont\bfseries\sffamily}
\renewcommand*{\parttitlefont}{\normalfont\fontsize{54pt}{54pt}\selectfont\bfseries\itshape\color{leadbeater}}
\renewcommand*{\printpartname}{%
  \partnamefont PART}

\makeatother

\chapterstyle{leadbeater}

\copypagestyle{leadbeater}{headings}

\makeoddhead{leadbeater}{\small\sffamily\color{leadbeater}\rightmark}{}
            {\color{leadbeater}\sffamily\bfseries\thepage}

\makeevenhead{leadbeater}{\small\sffamily\color{leadbeater}\leftmark}{}
            {\color{leadbeater}\sffamily\bfseries\thepage}

\usepackage{atbegshi}  % to kill the blank page after the titling page

\usepackage[font={small,it}]{caption}


% \olpath has to point to the location of the OLP main
% directory/folder.  We're compiling from subdirectory
% courses/phil379/, so the main directory is two
% levels up.

\newcommand{\olpath}{../../}

% load all the Open Logic definitions. This will also load the
% local definitions in open-logic-sample-config.sty

\input{\olpath/sty/open-logic.sty}

% we want all the problems deferred to the very end

\input{\olpath/sty/open-logic-defer.sty}

% colors for links

\hypersetup{
  allcolors=dkleadbeater,
%  citecolor=dkleadbeater,
%  urlcolor=dkleadbeater,
  pdftitle={Sets, Logic, Computation},
  pdfauthor = {Open Logic Project}
}

% load glossary entries

%\loadglsentries{include/glossary}

% end preamble

\input{includeonly}

\begin{document}

\AtBeginShipoutNext{\AtBeginShipoutNext{\AtBeginShipoutDiscard}}
% gets rid of empty page after title page

\begin{titlingpage}
  \pagecolor{leadbeater}
  \begin{textblock*}{14cm}(2cm,2cm)%
    \begin{raggedright}
      \fontsize{32pt}{34pt}\selectfont\bfseries\sffamily%
      Incompleteness and Computability\\
      \normalfont\fontsize{18pt}{0pt}\selectfont\bfseries\itshape%
      \rule{14cm}{5pt}\\[5pt]
      An Open Logic Text
    \end{raggedright}

%\vskip1cm\hskip-1cm\includegraphics[scale=.85]{illustrations/Cover}


    
\par\noindent
\vskip 15cm

\noindent
\includegraphics[width=1cm]{assets/cc.pdf}
\includegraphics[width=1cm]{assets/by.pdf}
\includegraphics[width=1cm]{assets/remix.pdf}
\normalfont\fontsize{16pt}{0pt}\selectfont\bfseries\sffamily%
\hfill Winter 2017% \textit{bis}
\end{textblock*}
\
\end{titlingpage}

\nopagecolor



% Now load the actual text

% phil379.tex
%
% driver file phil379.tex to produce text

\preto\OLEndChapterHook{\IfFileExists{include/summary-\thechapter}
        {\section*{Summary}\addcontentsline{toc}{section}{Summary}
        \let\emph\textbf\input{include/summary-\thechapter}\let\emph\textit}{}}

\problemsperchapter


\frontmatter

% bastard title

\pagestyle{empty}

\vspace*{100pt}

\begin{raggedleft}

{\fontsize{24pt}{24pt}\selectfont\bfseries\sffamily%
Incompleteness and
Computability}

\end{raggedleft}


\newpage

% editorial board

\vspace*{100pt}

{\bfseries\itshape The Open Logic Project}

\bigskip

\textbf{\color{leadbeater}Instigator}

\medskip

Richard Zach, \emph{University of Calgary}

\bigskip

\textbf{\color{leadbeater}Editorial Board}

\medskip

Aldo Antonelli,$^\dagger$ \emph{University of California, Davis}

Andrew Arana, \emph{Universit\'e Paris I Panth\'enon--Sorbonne}

Jeremy Avigad, \emph{Carnegie Mellon University}

Walter Dean, \emph{University of Warwick}

Gillian Russell, \emph{University of North Carolina}

Nicole Wyatt, \emph{University of Calgary}

Audrey Yap, \emph{University of Victoria}

\bigskip

\textbf{\color{leadbeater}Contributors}

\medskip

Samara Burns, \emph{University of Calgary}

Dana H\"agg, \emph{University of Calgary}

\newpage

% title

\vspace*{100pt}

\begin{raggedleft}

{\fontsize{24pt}{24pt}\selectfont\bfseries\sffamily%
  Incompleteness and Computability}

\smallskip

{\fontsize{18pt}{18pt}\selectfont\bfseries\itshape An Open Logic Text}

\vspace{100pt}

\fontsize{14pt}{14pt}\selectfont Remixed by Richard Zach

\vfill

\textsc{Winter 2017} %\textit{bis}

\end{raggedleft}


\newpage

% credits

\noindent
The Open Logic Project would like to acknowledge the generous support
of the \href{http://arts.ucalgary.ca/}{Faculty of Arts} and
the \href{http://www.ucalgary.ca/taylorinstitute/}{Taylor Institute of
Teaching and Learning} of the University of Calgary,
and the \href{http://albertaoer.com}{Alberta Open Educational Resources
(ABOER) Initiative}, which is made possible through an investment from
the Alberta government.



\bigskip

\noindent\includegraphics[height=1cm]{assets/ucarts-color}

\medskip

\noindent\includegraphics[height=1cm]{assets/ti-color}

\medskip

\noindent\includegraphics[width=4cm]{assets/aboer-color}

\vfill

\noindent Typeset in Baskervald X and Universalis ADF Standard by
\LaTeX.

\vfill

% oluselicense generates a license mark that a) licenses the result
% under a CC-BY licence and b) acknowledges the original source (the
% OLP).  Acknowledgment of the source is a requirement under the
% conditions of the CC-BY license used by the OLP, but you are not
% required to license the product itself under CC-BY.

\renewcommand{\ollicensefont}{\fontsize{9pt}{12pt}\selectfont}

\noindent
\oluselicense
% Title of this version of the OLT with link to source
{\href{https://github.com/rzach/phil479}{\textit{Incompleteness and
Computability}}}
% Author of this version
{\href{http://richardzach.org/}{Richard Zach}}

\newpage
\pagestyle{leadbeater}
\currentpdfbookmark{Table of Contents}{name}
\tableofcontents*

%\input{include/preface}



\pagestyle{leadbeater}

\mainmatter
\olimport*[incompleteness/introduction]{introduction}



\olimport*[computability/recursive-functions]{recursive-functions}



\olimport*[incompleteness/arithmetization-syntax]{arithmetization-syntax}



\olimport*[incompleteness/representability-in-q]{representability-in-q}



\olimport*[incompleteness/incompleteness-provability]{incompleteness-provability}



\appendix

\let\intro\comment
\let\endintro\endcomment

%\olimport*[first-order-logic]{first-order-logic}
% Part: first-order-logic
% Chapter: syntax-and-semantics

\chapter{First-order Logic}

\olimport*[first-order-logic/syntax-and-semantics]{first-order-languages}

\olimport*[first-order-logic/syntax-and-semantics]{terms-formulas}

\olimport*[first-order-logic/syntax-and-semantics]{free-vars-sentences}

\olimport*[first-order-logic/syntax-and-semantics]{substitution}

\olimport*[first-order-logic/syntax-and-semantics]{structures}

\olimport*[first-order-logic/syntax-and-semantics]{satisfaction}

\olimport*[first-order-logic/syntax-and-semantics]{semantic-notions}

\section{Theories}

\begin{defn}
A set of !!{sentence}s~$\Gamma$ is \emph{closed} iff, whenever
$\Gamma \Entails !A$ then $!A \in \Gamma$.  The \emph{closure} of a set
of !!{sentence}s~$\Gamma$ is $\Setabs{!A}{\Gamma \Entails !A}$.

We say that~$\Gamma$ is \emph{axiomatized by} a set of
sentences~$\Delta$ if $\Gamma$ is the closure of~$\Delta$
\end{defn}

\begin{ex}
The theory of strict linear orders in the language~$\Lang L_<$ is
axiomatized by the set
\begin{align*}
& \lforall[x][\lnot x < x], \\
& \lforall[x][\lforall[y][((x < y \lor y <
    x) \lor x = y)]], \\
& \lforall[x][\lforall[y][\lforall[z][((x < y
      \land y < z) \lif x < z)]]]
\end{align*}
It completely captures the intended !!{structure}s: every strict
linear order is a model of this axiom system, and vice versa, if $R$
is a linear order on a set $X$, then the structure $\Struct M$ with
$\Domain M = X$ and $\Assign{<}{M} = R$ is a model of this theory.
\end{ex}

\begin{ex}
The theory of groups in the language $\Obj 1$ (!!{constant}), $\cdot$
(two-place !!{function}) is axiomatized by
\begin{align*}
& \lforall[x][(x \cdot \Obj 1) = x]\\
& \lforall[x][\lforall[y][\lforall[z][\eq[(x \cdot (y \cdot z))][((x
          \cdot y) \cdot z)]]]]\\
& \lforall[x][\lexists[y][(x \cdot y) = \Obj 1]]
\end{align*}
\end{ex}

\begin{ex}
The theory of Peano arithmetic is axiomatized by the following
sentences in the language of arithmetic~$\Lang L_A$.
\begin{align*}
& \lnot\lexists[x][x' = \Obj 0]\\
& \lforall[x][\lforall[y][(x' = y' \lif x = y)]]\\
& \lforall[x][\lforall[y][(x < y \liff \lexists[z][(x + z' = y)])]]\\
& \lforall[x][\eq[(x + \Obj 0)][x]]\\
& \lforall[x][\lforall[y][\eq[(x + y')][(x + y)']]]\\
& \lforall[x][\eq[(x \times \Obj 0)][\Obj 0]]\\
& \lforall[x][\lforall[y][\eq[(x \times y')][((x \times y) + x)]]]\\
\intertext{plus all sentences of the form}
& (!A(\Obj 0) \land \lforall[x][(!A(x) \lif !A(x'))]) \lif \lforall[x][!A(x)]
\end{align*}
Since there are infinitely many sentences of the latter form, this
axiom system is infinite.  The latter form is called the
\emph{induction schema}. (Actually, the induction schema is a bit more
complicated than we let on here.)

The third axiom is an \emph{explicit definition} of~$<$.
\end{ex}


\OLEndChapterHook

\chapter{Natural Deduction}

\olimport*[first-order-logic/natural-deduction]{introduction}

\olimport*[first-order-logic/natural-deduction]{rules-and-proofs}

\olimport*[first-order-logic/natural-deduction]{proving-things}

\olimport*[first-order-logic/natural-deduction]{identity}

\olimport*[first-order-logic/natural-deduction]{proof-theoretic-notions}

\OLEndChapterHook




\stopproblems
\def\ifproblems#1{}

\def\captionnamefont{\sffamily\color{leadbeater}}
\def\captiontitlefont{\sffamily\color{leadbeater}}

\def\figurename{Fig.}
\setlength{\olphotowidth}{.45\textwidth}

\chapter{Biographies}

\olimport*[history/biographies]{alonzo-church}

\olimport*[history/biographies]{kurt-goedel}

\olimport*[history/biographies]{rozsa-peter}

\olimport*[history/biographies]{julia-robinson}

\olimport*[history/biographies]{alfred-tarski}

\backmatter

\clearpage
\photocredits

\bibliographystyle{\olpath/bib/natbib-oup}
\bibliography{\olpath/bib/open-logic.bib}

\olimport*{\olpath/content/open-logic-about}



\end{document}

